\section{Conclusions}

This paper presents the first theoretical analysis of bias in outcome disparity assessments using a probabilistic proxy model for the unobserved protected class.
In Theorem \ref{thm: demographic-est-race-bias}, we derived the bias of a thresholded estimator (Definition \ref{def:thresholded_estimator}) that has been described in the literature.
We also gave sufficient conditions in Corollary \ref{corollary: demographic-est-race} to understand when this methodology is biased,
and to what extent.
Our theoretical analysis is valid whenever a proxy model is used with a thresholded estimator to impute protected class membership,
and is thus consistent with previous studies that had observed an overestimation bias of the thresholded estimator based on BISG.
When applied to the public HMDA mortgage dataset with geolocation as the sole proxy for race membership (Section \ref{sec:mortgage}),
we found that the estimation bias of the thresholded estimator depends on the complex interaction of multiple different biases,
producing a strong sensitivity to the precise value of the threshold used.
Further studies will be needed to demonstrate the robustness of this numerical finding,
particularly when using other proxy models and other measures of fairness.
%Furthermore, the relative contribution of different bias sources is very sensitive to the threshold used in protected groups classification, which makes it very hard to assess the ultimate estimation bias.
Nevertheless, our work signals caution for choosing the \textit{ad hoc} thresholds which are used in practice,
as without the ground truth labels, we are unable to determine the optimal choice of threshold. 

%Our analysis of the thresholded estimator has revealed some theoretical issues with using hard thresholds for imputing the membership of a sample.
To alleviate the theoretical challenges of the thresholded estimator, we also proposed a weighted estimator (Definition 2.5),
which propagated the uncertainty resulting from the probabilistic proxy onto the final estimand.
We found that the estimation bias of this weighted estimator has only one bias source,
arising from the loan approval discrepancy between difference races within the same geolocation,
which led to an overall underestimate.
%. Overall, it is much more transparent when the weighted estimator is biased and unbiased compared to the thresholded estimator. 
As the behavior of this estimator's bias is simpler to reason about, 
we believe that the weighted estimator may be a useful new method to incorporate into outcome disparity evaluations when proxy models are used,
%  for the unobserved protected class,
especially if the sign of estimation bias can be determined with external knowledge.


% In the mortgage dataset where geolocation is used in the proxy model for race, one bias source results from the fact that geolocation encodes disparities in socioeconomic status, which are also related to creditworthiness, and this bias source contributes to the overestimation bias. The other bias source results from the fact that often the advantaged class has higher average loan approval rate than the disadvantaged class within the same geolocation, and this this bias source contribtues to the underestimation bias. 



% We presented herein the first theoretical analysis of the bias in outcome disparity assessments when a proxy model for the unknown protected attribute is employed.
% In Theorem \ref{thm: demographic-est-race-bias}, we stated sufficient conditions for the thresholded estimator (Definition \ref{def:thresholded_estimator}) to produce overestimates.
% We now have sufficient conditions in Corollary \ref{corollary: demographic-est-race} for understanding the empirical observations of when overestimation occurs \cite{zhang2016assessing,Baines:2014aa}.
% Intuitively, the geolocations used in probabilistic proxies such as BISG encode disparities in socioeconomic status, which are also related to creditworthiness.
% Thus, the estimated race labels end up correlated with the outcome being assessed.

% The analysis of the thresholded estimator revealed a source of bias due to the hard threshold needed to classify the membership of a sample; relaxing the hard threshold led naturally to the weighted estimator (Definition \ref{def:weighted_estimator}), which takes into account the inherent soft classification resulting from the use of a probabilistic proxy. We see empirically in numerical simulations that the weighted estimator tends to underestimate the amount of demographic disparity, whereas the thresholded estimator tends to overestimate. While we have not proven that the two estimators rigorously bound the true demographic disparity, we can nonetheless use the two estimators together to provide a more refined estimate for demographic disparity.

% Finally, we provide in Corollary \ref{corollary: demographic-est-race} and Proposition \ref{prop: relation-weighted-thresholding} some sufficient conditions under which the estimators for demographic disparity will overestimate or underestimate. The conditions set out in Corollary \ref{corollary: demographic-est-race} can be interpreted intuitively and can be used to analyze both toy examples and numerical experiments, even though the conditions may not be satisfied exactly.

%\section*{Disclaimer}
%
%None of the authors of this paper are lawyers.
%This paper does not constitute legal advice.
%The positions herein are presented for the purpose of academic research discussions and do not necessarily reflect the views of Capital One.
%
%\section*{Acknowledgments}
%
%This work was done when X.M.\ was at Capital One.
%We thank Elijah Alper, Brian Mills, Brian Larkin, Yu Huang, and Andrew Buemi for helpful discussions.
